\section{讨论与研究议程}
\subsection{集成、治理与隐私}

机制组合同时需要针对具体主张的策略与依赖分析。架构中的共同原因仍会产生影响：增加 UUV 数量，并不能消除它们对同一网关、时钟或运营方的依赖。

联盟设计必须规定准入、合约/模型批准、紧急权限、争议、成员退出和审计责任。海事采用研究与 TradeLens 经验说明，参与和激励为何需要纳入这项分析\citep{zhou2020key,farah2024survey,maersk2022tradelens}。示例中的背书规则保留共同控制，同时通过必需签名方产生可测量的可用性成本。同一所有者控制的更多副本，不会形成更多独立治理域。

跨行业研究区分可追溯性收益与实现这些收益所需的组织变革\citep{kshetri20181,lnes2017blockchain}。联盟协议分配留存、披露、维护和事件响应职责，包括成员退出后的证据可用性，以及密钥轮换后的历史签名处理。比较账本与验证日志时，必须匹配这些治理规定。

隐私取决于载荷加密之外的因素。稳定标识符、路径托管、背书时序和模型使用记录，可能泄露无人平台活动与组织关系\citep{reyna2018blockchain,kshetri2017can}。因此，可审计性需要同时配备数据最小化、选择性披露、留存规则和受控身份解析。撤销依据策略使未来依赖失效；它不会删除已保留的承诺。人工覆盖决定是一项另行授权并记录的转换，应包含证据、理由和后果。

\subsection{可证伪的研究命题}\label{sec:propositions}

以下命题定义未来的比较测试。每个实验预先规定任务分布、最小有用效应、主要结果和分析规则。零效应和不利结果能够识别附加结构不产生收益的条件。

\textbf{P1（主张组织）。}检验保障向量能否在匹配的错误隔离与时延预算下，相对于调优后的标量策略，降低留出测试中的不安全接受。两者接收完整证据、相同行动集合和相等调参预算；标量包含授权/时效性门控与交互项。若比较方法表现相同或更好，则命题被否证，表~\ref{tab:trace} 中带否决条件的标量就是这样的例子。

\textbf{P2（依赖建模）。}在预先规定的共享源误差和合谋率下，检验含共同原因变量的模型，相对于同等调优的独立性模型，是否改善适当评分规则和任务损失。纳入正确、不完整和错误谱系条件。当模型设定错误或采集成本抵消收益，或观测足够独立以至于不存在有用改进时，命题不成立。

\textbf{P3（诊断与未知状态）。}在退化、攻击和证据中断混合发生时，比较保留非恶意、对抗性和未决假设的模型与二元评估模型。两者都可以选择相同行动，包括弃权。分别消融诊断与弃权功能，并测量不安全接受、错误隔离、佐证延迟和可用性损失带来的损失。若增加诊断状态导致损失上升，或未带来有用改进，则命题不成立。

\textbf{P4（离线权限）。}改变委托范围、到期时间、时钟不确定性、分区持续时间和撤销时点。与仅允许新获取全局授权，以及不受限制的离线自主运行相比，检验有界委托是否能够在截止时间前完成的任务与违反已生效全局策略的行动之间，提供有用的权衡。分别报告本地获授权行动与事后发现冲突的行动。若某些运行条件下，有界策略在这两个结果上都被其他方法支配，则包括在命题失败的情形中。

\textbf{P5（条件性的账本价值）。}在匹配成员资格、故障、规则、留存和工作负载的条件下，比较已实现的账本与独立见证日志服务。改变 $n$、有响应的法定人数、分区持续时间 $D$、事件率和积压量。对于截止时间 $\tau$，使用在负载 $\lambda$ 下实测的恢复接触后延迟 $L$，估计 $D+L(n,\lambda)\leq\tau$。在固定字节/能耗预算下，检验及时可验证决策或恢复完整性的改进。若安全性相同但成本更高、背书方可用性更低，或相对于日志没有改进，则该运行条件下的优势被否证。第~\ref{sec:budget} 节的 58/59 s 边界使用的是假定延迟。

\textbf{P6（评估覆盖）。}在相同候选集合与相等评估预算下，分别仅使用事务/网络指标，或在其基础上加入观测、端点、治理和任务结果，来选择配置。按固定选择规则比较留出测试中的任务损失。若计入额外采集与解释成本后，更全面的评估未改善选择，则命题不成立。

\subsection{保障论证与海事要求}\label{sec:assurance-case}

OAL 提供保障论证的内容，已有表示法提供其结构。GSN 表示主张、策略、情境、假设和证据，SACM 则提供结构化保障论证元模型\citep{gsn2021,sacm2023}。UL 4600 提供自主产品的安全论证要求，而非图形表示法\citep{ul4600}。OAL 可以用对象/版本、保障层、生命周期阶段、事件区间和证据类型标注它们的主张。符合性评估仍与适用标准的要求相关联。

航线重规划的适用性需要可信观测、受保护的派生过程、有效权限、当前策略和可接受的不确定性。论证记录观测、签发、接收和到期时间，包括缓存权限。输入更正或密钥撤销会触发派生后代审查，同时保留具有独立支撑的产品。验证应测量证据缺口、对植入的无效论证的检测、审查者错误，以及识别受影响产品所需时间。

海事要求增加了特定应用边界。IMO 的海上自主水面船舶（MASS）试验临时指南和监管范围界定工作，提供试验与监管背景\citep{imo2019trials,imo2021scoping}。非强制性 MASS 规则于 2026 年 5 月通过，并于 2026 年 7 月 1 日生效\citep{imo2026mass}。DNV-CG-0264 面向自主及遥控船舶\citep{dnv2024autonomous}。保障论证必须明确这些文件所适用的船舶和作业、责任主管机构、运行设计条件、人员/远程控制职责，以及安全回退证据。其海事适用范围不能整体移用于 UAV 或水下作业。账本记录可以追踪批准或运行限制；实质性安全证据和适用审批流程仍需分别处理。

\subsection{迈向跨域无人系统信任配置规范}

提出的\CDUTP{}是一项关于版本化主张与证据引用的互操作研究议程。它可以复用 RATS 角色、W3C PROV 关系、零信任策略概念和安全更新清单\citep{birkholz2023remote,w3c2013provo,rose2020zero,ietf2020suit}。UAS 与海事服务规范提供通信背景\citep{threegpp22125,threegpp22119}。一个配置规范实例记录对象/类型、限定范围的主张、签发方和依赖角色、证据引用、事件与接收时间、到期时间、不确定性、允许用途、策略/模型/模式版本、披露条件、依赖关系和撤销绑定。

第~\ref{sec:pattern} 节的紧凑编码是一种示例传输配置。互操作性需要针对具体版本的语义映射与测试，包括有界时钟、密钥轮换、证据缺失和断连签发。随后可在互操作联测与多组织演练中比较独立实现。核心研究问题是：共同语义能否在可接受的通信与治理成本下，减少决策和审计错误。
